\documentclass[10pt,a4paper]{article}
\usepackage[margin=2.2cm]{geometry}
\usepackage{fontspec}
\setmonofont{DejaVu Sans Mono}[Scale=0.84]
\newfontfamily\unifont{DejaVu Serif}[Scale=0.9]
\usepackage{amsmath,amssymb,mathtools}
\setcounter{MaxMatrixCols}{20}
\usepackage{longtable,booktabs,array}
\usepackage{fancyvrb}
\usepackage[dvipsnames]{xcolor}
\usepackage[colorlinks=true,linkcolor=NavyBlue,urlcolor=NavyBlue]{hyperref}
\usepackage{microtype}
\usepackage{needspace}
\usepackage{tocloft}\setlength{\cftbeforesecskip}{0.45em}
\setlength{\parindent}{0pt}\setlength{\parskip}{0.55em}
\setlength{\emergencystretch}{8em}
\makeatletter
\renewcommand\section{\@startsection{section}{1}{\z@}{-3.5ex \@plus -1ex \@minus -.2ex}{2.3ex \@plus.2ex}{\normalfont\Large\bfseries\raggedright}}
\renewcommand\subsection{\@startsection{subsection}{2}{\z@}{-3.25ex\@plus -1ex \@minus -.2ex}{1.5ex \@plus .2ex}{\normalfont\large\bfseries\raggedright}}
\renewcommand\subsubsection{\@startsection{subsubsection}{3}{\z@}{-3.25ex\@plus -1ex \@minus -.2ex}{1.5ex \@plus .2ex}{\normalfont\normalsize\bfseries\raggedright}}
\makeatother
\renewcommand{\arraystretch}{1.15}
\allowdisplaybreaks
\begin{document}
\title{7+1 repository after the clean-up of 28 October 2026: symbols and terms\\[0.4em]\large A dictionary of every symbol and term of the document on the state of the repository at commit 102f399 and of the spreadsheet 7plus1\_status\_2026-10-28.xlsx, with the defining equations}
\author{Repository: github.com/once-ere/7plus1 — prepared for Nash\\ Commit described: 102f399 of 28 October 2026, 13:15 PDT (branch main)\\ Document generated by tools/make\_docs.py on 2026-10-28}
\date{}
\maketitle
\begin{abstract}
For each symbol or term this dictionary gives its definition and its meaning, grouped by topic: the 7+1 spacetime, the repository and its commits, the documents and TeX’s measure of a loose line, the kinds of rewritten paragraphs, the checks, the spreadsheet and its functions, and every symbol of the formulas of the rewritten paragraphs. Longer definitions follow each table as numbered displayed equations. Section 10 lists the file and program names, section 11 the symbols that look alike. Every equation needed to define a symbol is written in this document.
\end{abstract}
\tableofcontents
\Needspace{12\baselineskip}\section*{1. Conventions}\addcontentsline{toc}{section}{1. Conventions}
\begin{itemize}
\item Coordinates $X=(x_0,x_1,\dots,x_7)$: $x_0$ is time, the direction of the octonion unit $e_0=1$; $x_1,x_2,x_3$ ordinary space; $x_4,\dots,x_7$ the four auxiliary space-like dimensions. Indices: $A,B=0,\dots,7$; $\mu,\nu,\rho,\sigma,\lambda=0,\dots,7$ for coordinate indices; $i,j=1,\dots,7$ for the space directions; repeated indices are summed unless stated otherwise.
\item Units: $\hbar=c=1$; masses and energies in GeV.
\item Matrices: $1_n$ is the $n\times n$ identity, $M^{\mathsf T}$ the transpose, $M^\dagger$ the conjugate transpose, $M^*$ the complex conjugate (section 9.14 writes them out); $\delta_{ij}$ is the Kronecker delta (1 for $i=j$, 0 otherwise).
\item Christoffel symbols are written $\{{}^{\lambda}_{\mu\nu}\}$, never $\Gamma$: in this project $\Gamma$ always denotes a Dirac matrix.
\end{itemize}
\textbf{The signature: $\{7,1,0\}$ and $\mathrm{diag}(+1,-1,\dots,-1)$.} The 7+1 spacetime has one time-like dimension ($x_0$, the direction of the octonion unit $e_0=1$) and seven space-like dimensions ($x_1,x_2,x_3$ ordinary space, $x_4,\dots,x_7$ auxiliary space). Two different objects describe this, and neither may be replaced by the other. (i) The tensor package xAct, used in section 12 of 7+1.nb exactly as in the author’s notebook Einstein-Lovelock-7+1-Nash.nb, declares the metric with \texttt{DefMetric[\{7, 1, 0\}, g71[-a, -b], CD, …]}: the list $\{7,1,0\}$ declares seven space-like dimensions, one time-like dimension and no degenerate dimension — the numbers of entries $+1$ and $-1$ of $\eta_{71}=\mathrm{diag}(-1,+1,\dots,+1)$ (\texttt{η71}), the same spacetime written in the sign convention of 4+4.nb (space-like $=+1$, time-like $=-1$). (ii) The matrix $\eta=-\eta_{71}=\mathrm{diag}(+1,-1,\dots,-1)$ (\texttt{η17}) is simply the squares of the octonion basis: $e_0^2=+1$ for the unit $e_0=1$ (time) and $e_i^2=-1$ for the seven imaginary units $e_1,\dots,e_7$ (space); more generally it is the real part of the product of two basis units, and it gives the quadratic form of the octonions:
\begin{gather*}
\eta_{AB}=\mathrm{Re}(e_Ae_B),\qquad e_0^2=+1,\qquad e_i^2=-1\ \ (i=1,\dots,7),\\ \mathrm{Re}(x\,x)=x_0^2-x_1^2-\cdots-x_7^2=\eta_{AB}\,x^Ax^B\qquad (x=x^Ae_A\in\mathbb O).
\end{gather*}
$\eta$ gives the signs of the Clifford relation $\Gamma_A\Gamma_B+\Gamma_B\Gamma_A=2\,\eta_{AB}\,1_{16}$ of the real 16×16 Dirac matrices: time squares to $+1_{16}$ and each of the seven space directions to $-1_{16}$, exactly like the octonion units (in the notation $Cl(p,q)$, with $p$ generators squaring to $+1$ and $q$ to $-1$, the algebra $Cl(1,7)\cong\mathbb R(16)$). It also gives the signs of the metric components: $g_{00}>0$ along the time-like direction, negative squares along the space-like directions. Every component computed by 7+1.nb follows from the component matrix of the metric; the declared list fixes for xAct the sign of the determinant of the metric, \texttt{SignDetOfMetric[g71]} $=-1$, which is the same in both sign conventions, since $\det \eta=\det \eta_{71}=-1$. 7+1.nb checks \texttt{SignatureOfMetric[g71]} $=\{7,1,0\}$, \texttt{SignDetOfMetric[g71]} $=-1$, $\mathrm{Re}(e_Ae_B)=\eta_{AB}$ and $\mathrm{Re}(x\,x)=\eta_{AB}x^Ax^B$.

\Needspace{12\baselineskip}\section*{2. The repository and its commits}\addcontentsline{toc}{section}{2. The repository and its commits}
{\footnotesize
\begin{longtable}{>{\raggedright\arraybackslash}p{\dimexpr 0.2400\linewidth-2\tabcolsep\relax}>{\raggedright\arraybackslash}p{\dimexpr 0.3600\linewidth-2\tabcolsep\relax}>{\raggedright\arraybackslash}p{\dimexpr 0.4000\linewidth-2\tabcolsep\relax}}
\toprule
\textbf{symbol or term} & \textbf{definition} & \textbf{meaning} \\
\midrule\endfirsthead
\toprule
\textbf{symbol or term} & \textbf{definition} & \textbf{meaning} \\
\midrule\endhead
\endfoot
\bottomrule\endlastfoot
repository & github.com/once-ere/7plus1 & the Git repository of the 7+1 project, on GitHub \\
branch \texttt{main} & the line of commits that is published & the commits of the repository \\
commit & a recorded state of all files & one step of the history; its parent is the commit before \\
commit name (\texttt{102f399}) & the first 7 hexadecimal digits of the 40-digit SHA-1 object name & short name; the full name of 102f399 is \texttt{102f399d\allowbreak{}53720ba9\allowbreak{}b6bc05c6\allowbreak{}dbf8d70c\allowbreak{}68c9ac30} \\
commit described & 102f399 & the state of the repository that the spreadsheet records \\
$F_c$, $A_c$, $R_c$ & (2.1) & files changed, lines added, lines removed by the commit $c$ \\
$F$, $A$, $R$ & (2.2) & the totals of the four commits \\
file change & one file changed by one commit & a file changed by two commits is counted twice \\
binary file & a file that Git does not compare line by line (a PDF) & counts as changed, adds and removes no lines \\
fresh clone & \texttt{git clone https:/\allowbreak{}/\allowbreak{}github.com/\allowbreak{}once-ere/\allowbreak{}7plus1.git} into an empty folder & a copy that contains only what is published \\
tracked file & a file listed by \texttt{git ls-files} & a file that is part of the repository \\
\texttt{.gitignore} & a file of patterns & files that Git does not track; this one is not published itself \\
\texttt{.git/\allowbreak{}info/\allowbreak{}exclude} & a file of patterns in one working folder & files that Git does not track there; never published \\
private file & a file of the author in the working folder, not tracked & not on GitHub \\
SHA-256 & a 256-bit cryptographic hash, written as 64 hexadecimal digits & the fingerprint of the contents of a file \\
PDT & UTC−7 & Pacific Daylight Time, the time zone of the author’s computer \\*
\end{longtable}}
\begin{align*}
(F_c,\ A_c,\ R_c)&=\text{the numbers of }\texttt{git show --shortstat}\ \text{for the commit }c,&&(2.1) \\
F=\sum_cF_c,\qquad A&=\sum_cA_c,\qquad R=\sum_cR_c\qquad(c=\texttt{0654dcd},\ \texttt{54862f0},\ \texttt{93c500f},\ \texttt{102f399}).&&(2.2)
\end{align*}
For the four commits of 28 October 2026: $F=122+12+44+28=206$, $A=24996+66+210+374=25646$, $R=4890+19+93+355=5357$.

\Needspace{12\baselineskip}\section*{3. Documents and typesetting}\addcontentsline{toc}{section}{3. Documents and typesetting}
{\footnotesize
\begin{longtable}{>{\raggedright\arraybackslash}p{\dimexpr 0.2400\linewidth-2\tabcolsep\relax}>{\raggedright\arraybackslash}p{\dimexpr 0.3600\linewidth-2\tabcolsep\relax}>{\raggedright\arraybackslash}p{\dimexpr 0.4000\linewidth-2\tabcolsep\relax}}
\toprule
\textbf{symbol or term} & \textbf{definition} & \textbf{meaning} \\
\midrule\endfirsthead
\toprule
\textbf{symbol or term} & \textbf{definition} & \textbf{meaning} \\
\midrule\endhead
\endfoot
\bottomrule\endlastfoot
document & one documentation file, as \texttt{.md}, \texttt{.tex} and \texttt{.pdf} & Markdown for GitHub, LaTeX source, compiled PDF \\
kind & Main, Symbols or Provenance & a document, its symbol dictionary, its provenance record \\
covers & sections of 7+1.nb & what the document describes \\
pages & the number reported by \texttt{pdfinfo} & the length of the PDF \\
$w$ & the width of a line & the box that TeX fills \\
$w_0$ & natural width & the width of the material of the line with every space at its natural size \\
$s$, $z$ & total stretchability, total shrinkability & by how much the spaces of the line may grow or shrink \\
$t$ & $w-w_0$ (when positive) & the stretch that the line needs \\
$b$ & (3.1) & badness: how much a line is stretched, from 0 (natural) to 10000 (infinitely bad) \\
\texttt{\textbackslash{}\allowbreak{}hbadness} & 1000 (set by LaTeX) & the badness above which TeX reports a box \\
\texttt{\textbackslash{}\allowbreak{}hfuzz} & 0.1 pt (set by LaTeX) & the excess width that TeX tolerates silently \\
loose line & (3.2); the message \texttt{Underfull \textbackslash{}\allowbreak{}hbox} & a line whose spaces stretch too much \\
overflowing line & (3.3); the message \texttt{Overfull \textbackslash{}\allowbreak{}hbox} & a line wider than the line width \\
LaTeX error & a line of the log that begins with “!” & a fault in the LaTeX source \\
layout problem & a loose or overflowing line, an error, or a missing PDF & what the rebuild must not show \\
\texttt{latexmk -xelatex} & the program that runs XeLaTeX until the cross-references settle & compiles each \texttt{.tex} file inside build/docs \\
\texttt{pdftotext} & the text of a PDF & used to compare two PDFs \\
first build & the first full build of the documents on 28 October 2026 & before commit 0654dcd: 40 loose lines \\
changed by a commit & the commit changed the \texttt{.md}, \texttt{.tex} or \texttt{.pdf} file of the document &  \\*
\end{longtable}}
\begin{align*}
b&=\min\Big(10000,\ 100\Big(\frac ts\Big)^3\Big)\ \ (s>0),\qquad b=10000\ \ (s=0,\ t>0),&&(3.1) \\
\text{loose}&\iff b>1000\iff\frac ts>10^{1/3}\approx2.154,&&(3.2) \\
\text{overflowing}&\iff w_0-z-w>0.1\ \text{pt}.&&(3.3)
\end{align*}
TeX rounds $b$ to an integer. The number of loose (overflowing) lines of a document is the number of messages \texttt{Underfull \textbackslash{}\allowbreak{}hbox} (\texttt{Overfull \textbackslash{}\allowbreak{}hbox}) in the log of its last LaTeX run; each message is one box, a line of a paragraph or of a table cell.

\Needspace{12\baselineskip}\section*{4. The kinds of rewritten paragraphs}\addcontentsline{toc}{section}{4. The kinds of rewritten paragraphs}
{\footnotesize
\begin{longtable}{>{\raggedright\arraybackslash}p{\dimexpr 0.1600\linewidth-2\tabcolsep\relax}>{\raggedright\arraybackslash}p{\dimexpr 0.4000\linewidth-2\tabcolsep\relax}>{\raggedright\arraybackslash}p{\dimexpr 0.4400\linewidth-2\tabcolsep\relax}}
\toprule
\textbf{“How”} & \textbf{definition} & \textbf{example} \\
\midrule\endfirsthead
\toprule
\textbf{“How”} & \textbf{definition} & \textbf{example} \\
\midrule\endhead
\endfoot
\bottomrule\endlastfoot
Line breaking & a change of the typesetting rules of all documents & code may break after \texttt{/\allowbreak{}}, \texttt{\_\allowbreak{}} and \texttt{\textbackslash{}\allowbreak{}}, a long hexadecimal string every 8 characters; headings are set ragged right (not stretched to the full width) \\
Display & a long inline formula moved into a displayed equation & the definitions of $\Theta^K$, $\sigma_r$ and $\Theta_K$ (§9.1) \\
Code block & a long piece of code moved into a block of its own & the xAct definitions (§9.3) \\
Reworded & the sentence rewritten so that TeX can break its lines evenly & the electroweak scale (§9.13) \\
Correction & a correction of content & the check counts 50 and 38 (before, 49 and 37) \\*
\end{longtable}}
A paragraph that is rewritten keeps its content. “Paragraphs rewritten” counts the rows of the kinds Display, Code block and Reworded.

\Needspace{12\baselineskip}\section*{5. The checks}\addcontentsline{toc}{section}{5. The checks}
{\footnotesize
\begin{longtable}{>{\raggedright\arraybackslash}p{\dimexpr 0.2400\linewidth-2\tabcolsep\relax}>{\raggedright\arraybackslash}p{\dimexpr 0.3600\linewidth-2\tabcolsep\relax}>{\raggedright\arraybackslash}p{\dimexpr 0.4000\linewidth-2\tabcolsep\relax}}
\toprule
\textbf{symbol or term} & \textbf{definition} & \textbf{meaning} \\
\midrule\endfirsthead
\toprule
\textbf{symbol or term} & \textbf{definition} & \textbf{meaning} \\
\midrule\endhead
\endfoot
\bottomrule\endlastfoot
check program & one of the twelve Python programs of tools/ (section 10) & recomputes results independently, or parses the notebook \\
exit code & the status with which a program ends & 0 means success \\
$\ell$ & a line of the output of a program &  \\
$N_{\rm True}(P)$, $N_{\rm False}(P)$ & (5.1) & the results reported True and False by the program $P$ \\
verdict & the last line \texttt{VERDICT: …} of the output, or, without it, “All printed results True” & the summary of a program \\
notebook check & a cell \texttt{check["name", test]} of 7+1.nb, or a check made by \texttt{frameChecks} & a test that returns True or False \\
$N_{\rm nb}$, $N_{\rm printed}$, $N_{\rm frame}$ & (5.2) & checks of the notebook; printed where they stand; made by \texttt{frameChecks} \\
\texttt{frameChecks} & a function of section 13.1 of 7+1.nb & nine checks of the frame matrices for each of four metrics \\
input, output cell & cells of the notebook & the Wolfram Language code and its results \\
message cell & a cell in which Mathematica reports a problem & must be 0 \\
fresh kernel & a new Wolfram Language kernel with no earlier definitions & the evaluation starts from nothing \\
evaluation time & the duration of the evaluation of all cells & 283.883 s for the evaluation that finished on 2026-10-28 08:57:47 PDT \\
rebuild & all documents written anew by \texttt{tools/\allowbreak{}make\_\allowbreak{}docs.py} in a fresh clone & compared with the committed files \\
identical byte for byte & the same sequence of bytes & the comparison of Markdown and LaTeX files \\
same text and page count & equal \texttt{pdftotext} output and equal \texttt{pdfinfo} page count & the comparison of PDFs, which also record their creation time \\*
\end{longtable}}
\begin{align*}
N_{\rm True}(P)&=\#\{\ell:\ \ell\ \text{ends in “: True” or in “(True)”}\},&&(5.1) \\
N_{\rm False}(P)&=\#\{\ell:\ \ell\ \text{ends in “: False” or in “(False)”}\}, \\
N_{\rm nb}&=N_{\rm printed}+N_{\rm frame},\qquad N_{\rm frame}=9\times4\qquad(479=443+36).&&(5.2)
\end{align*}
In (5.1) $\#$ counts the elements of a set, $\ell$ runs over the lines of the output of $P$, and spaces between the colon and the word, and after the word, are allowed.

\Needspace{12\baselineskip}\section*{6. The spreadsheet}\addcontentsline{toc}{section}{6. The spreadsheet}
{\footnotesize
\begin{longtable}{>{\raggedright\arraybackslash}p{\dimexpr 0.2400\linewidth-2\tabcolsep\relax}>{\raggedright\arraybackslash}p{\dimexpr 0.3600\linewidth-2\tabcolsep\relax}>{\raggedright\arraybackslash}p{\dimexpr 0.4000\linewidth-2\tabcolsep\relax}}
\toprule
\textbf{symbol or term} & \textbf{definition} & \textbf{meaning} \\
\midrule\endfirsthead
\toprule
\textbf{symbol or term} & \textbf{definition} & \textbf{meaning} \\
\midrule\endhead
\endfoot
\bottomrule\endlastfoot
sheet & Summary, Commits, Documents, Fixes, Checks, Private files & the six tables of the spreadsheet \\
cell \texttt{B5} & column B, row 5 & one field of a sheet \\
\texttt{Commits!E6} & the cell E6 of the sheet Commits & a reference to another sheet \\
\texttt{'Private files'!C2} & a sheet name with a space, quoted &  \\
range \texttt{A2:A5} & the rectangle of cells from A2 to A5 &  \\
$R$, $c$, $v(c)$ & a range, a cell, the value of the cell &  \\
formula & a cell whose content begins with \texttt{=} & its value is computed from other cells \\
\texttt{SUM}, \texttt{COUNTA}, \texttt{COUNTIF} & (6.1) & the functions of the spreadsheet \\
\texttt{+}, \texttt{\&} & (6.2) & addition; joining of texts \\
data file & tools/status\_2026-10-28.json & every number of the spreadsheet \\
openpyxl & a Python library that writes \texttt{.xlsx} files & writes the spreadsheet with its formulas \\
LibreOffice (headless) & the office suite, run without a window & opens and saves the file, which computes and stores the formula values \\
stored value & the value of a formula saved in the file & what a viewer shows before computing anything \\*
\end{longtable}}
\begin{align*}
\mathrm{SUM}(R)&=\sum_{c\in R,\ v(c)\ \text{a number}}v(c),\qquad \mathrm{COUNTA}(R)=\#\{c\in R:\ v(c)\ \text{not empty}\},&&(6.1) \\
\mathrm{COUNTIF}(R,x)&=\#\{c\in R:\ v(c)=x\}, \\
x+y&=\text{the sum of }x\text{ and }y,\qquad x\,\&\,y=\text{the text of }x\text{ followed by the text of }y .&&(6.2)
\end{align*}
In \texttt{COUNTIF}, texts are compared without regard to case.

\Needspace{12\baselineskip}\section*{7. The time line of the clean-up}\addcontentsline{toc}{section}{7. The time line of the clean-up}
{\footnotesize
\begin{longtable}{>{\raggedright\arraybackslash}p{\dimexpr 0.1800\linewidth-2\tabcolsep\relax}>{\raggedright\arraybackslash}p{\dimexpr 0.6200\linewidth-2\tabcolsep\relax}>{\raggedright\arraybackslash}p{\dimexpr 0.2000\linewidth-2\tabcolsep\relax}}
\toprule
\textbf{moment} & \textbf{definition} & \textbf{loose lines} \\
\midrule\endfirsthead
\toprule
\textbf{moment} & \textbf{definition} & \textbf{loose lines} \\
\midrule\endhead
\endfoot
\bottomrule\endlastfoot
first build & the first full build of the documents on 28 October 2026, before commit 0654dcd & 40 \\
after 0654dcd & the rebuild of a fresh clone of 0654dcd & 12 \\
now & the rebuild of a fresh clone of 102f399 & 0 \\*
\end{longtable}}
\Needspace{12\baselineskip}\section*{8. Numbers and constants}\addcontentsline{toc}{section}{8. Numbers and constants}
{\footnotesize
\begin{longtable}{>{\raggedright\arraybackslash}p{\dimexpr 0.2000\linewidth-2\tabcolsep\relax}>{\raggedright\arraybackslash}p{\dimexpr 0.3600\linewidth-2\tabcolsep\relax}>{\raggedright\arraybackslash}p{\dimexpr 0.4400\linewidth-2\tabcolsep\relax}}
\toprule
\textbf{symbol} & \textbf{value} & \textbf{meaning} \\
\midrule\endfirsthead
\toprule
\textbf{symbol} & \textbf{value} & \textbf{meaning} \\
\midrule\endhead
\endfoot
\bottomrule\endlastfoot
$G_F$ & $1.1663788\times10^{-5}$ GeV$^{-2}$ & Fermi constant \\
$v$ & 246.2196402 GeV & electroweak scale, $(\sqrt2G_F)^{-1/2}$ \\
$M_{\rm Pl}$ & $2.435\times10^{18}$ GeV & reduced Planck mass, $8\pi G_4=M_{\rm Pl}^{-2}$ \\
$p_1$, $q_1$, $u_1$ & $\frac{1+2\sqrt2}7$, $u_1p_1$, $-1+\frac1{\sqrt2}$ & Kasner constants, (9.3) \\
$10^{1/3}$ & 2.154 & the stretch ratio of badness 1000 \\*
\end{longtable}}
\Needspace{12\baselineskip}\section*{9. The symbols of the formulas of the rewritten paragraphs}\addcontentsline{toc}{section}{9. The symbols of the formulas of the rewritten paragraphs}
“Row $n$” is the $n$-th row of the sheet Fixes, that is, the $n$-th rewritten paragraph. Each subsection gives the formula of a row as its document writes it, then its symbols; “(row $n$)” in the column “definition” means that the formula at the top of the subsection defines the symbol.

\Needspace{8\baselineskip}\subsection*{9.1 The Clifford products (row 2)}\addcontentsline{toc}{subsection}{9.1 The Clifford products (row 2)}
\begin{equation*}
\Theta^K=i^{\sigma_r}\Gamma^{A_1}\cdots\Gamma^{A_r},\qquad \sigma_r=\frac{r(r-1)}2\bmod2,\qquad \Theta_K=\eta_{A_1A_1}\cdots\eta_{A_rA_r}\Theta^K .
\end{equation*}
{\footnotesize
\begin{longtable}{>{\raggedright\arraybackslash}p{\dimexpr 0.2400\linewidth-2\tabcolsep\relax}>{\raggedright\arraybackslash}p{\dimexpr 0.3600\linewidth-2\tabcolsep\relax}>{\raggedright\arraybackslash}p{\dimexpr 0.4000\linewidth-2\tabcolsep\relax}}
\toprule
\textbf{symbol or term} & \textbf{definition} & \textbf{meaning} \\
\midrule\endfirsthead
\toprule
\textbf{symbol or term} & \textbf{definition} & \textbf{meaning} \\
\midrule\endhead
\endfoot
\bottomrule\endlastfoot
$\Gamma_A=\Gamma^{(1,7)}_A$ & (9.1) & the eight real 16×16 Dirac matrices of 7+1.nb, $A=0,\dots,7$ \\
$\tau^O_i$ & real antisymmetric 8×8 matrices, (9.1) & the seven imaginary octonion units $e_1,\dots,e_7$ as matrices \\
$\eta_{AB}$ & $\mathrm{diag}(+1,-1,\dots,-1)$ & the squares of the octonion basis (section 1) \\
$\Gamma^A$ & $\eta^{AA}\Gamma_A$ (no sum) & Dirac matrix with the index raised: $\Gamma^0=\Gamma_0$, $\Gamma^i=-\Gamma_i$ \\
$K=(A_1<\dots<A_r)$ & an increasing list of indices & a multi-index; there are $\binom8r$ of rank $r$ and 256 in all \\
$r=|K|$ & $0,\dots,8$ & the rank of the multi-index \\
$\sigma_r$ & $\frac{r(r-1)}2\bmod2$ & 0 or 1: the phase $i^{\sigma_r}$ \\
$\Theta^K$ & $i^{\sigma_r}\Gamma^{A_1}\cdots\Gamma^{A_r}$ & the 256 Clifford products ($\Theta^{()}=1_{16}$) \\
$\Theta_K$ & $\eta_{A_1A_1}\cdots\eta_{A_rA_r}\Theta^K$ & index lowered; $\Theta_K\Theta^K=1_{16}$ \\
$i$ & $i^2=-1$ & the imaginary unit \\*
\end{longtable}}
\begin{align*}
\Gamma^{(1,7)}_0&=\begin{pmatrix}0&1_8\\1_8&0\end{pmatrix},\qquad \Gamma^{(1,7)}_i=\begin{pmatrix}\tau^O_i&0\\0&-\tau^O_i\end{pmatrix}\quad(i=1,\dots,7),&&(9.1) \\
\tau^O_i\tau^O_j+\tau^O_j\tau^O_i&=-2\delta_{ij}1_8,\qquad \Gamma_A\Gamma_B+\Gamma_B\Gamma_A=2\eta_{AB}1_{16}.
\end{align*}
\Needspace{8\baselineskip}\subsection*{9.2 The Grassmann algebra (row 3)}\addcontentsline{toc}{subsection}{9.2 The Grassmann algebra (row 3)}
\begin{align*}
\theta_K\theta_{K'}&=\mathrm{sign}(\pi)\,\theta_{\mathrm{sort}(K\cup K')}\qquad(\text{zero if }K\cap K'\neq\emptyset), \\
\frac{\partial^L(\theta_{k_1}\cdots\theta_{k_r})}{\partial\theta_{k_j}}&=(-1)^{j-1}\,\theta_{k_1}\cdots\widehat{\theta_{k_j}}\cdots\theta_{k_r} .
\end{align*}
{\footnotesize
\begin{longtable}{>{\raggedright\arraybackslash}p{\dimexpr 0.2400\linewidth-2\tabcolsep\relax}>{\raggedright\arraybackslash}p{\dimexpr 0.3600\linewidth-2\tabcolsep\relax}>{\raggedright\arraybackslash}p{\dimexpr 0.4000\linewidth-2\tabcolsep\relax}}
\toprule
\textbf{symbol or term} & \textbf{definition} & \textbf{meaning} \\
\midrule\endfirsthead
\toprule
\textbf{symbol or term} & \textbf{definition} & \textbf{meaning} \\
\midrule\endhead
\endfoot
\bottomrule\endlastfoot
$\theta_k$ & $\theta_k\theta_l=-\theta_l\theta_k$, so $\theta_k^2=0$ & the anticommuting generators \\
$K$, $K'$ & increasing lists of generator indices & index sets \\
$\theta_K$ & $\theta_{k_1}\cdots\theta_{k_r}$ for $K=(k_1<\dots<k_r)$ & the basis monomials \\
$\pi$ & the permutation that sorts $K$ followed by $K'$ &  \\
$\mathrm{sign}(\pi)$ & $+1$ for an even, $-1$ for an odd permutation &  \\
$\mathrm{sort}(K\cup K')$ & the union in increasing order &  \\
$\partial^L/\partial\theta_{k_j}$ & the left derivative & moves $\theta_{k_j}$ to the front, then removes it \\
$\widehat{\theta_{k_j}}$ & the factor left out &  \\*
\end{longtable}}
\Needspace{8\baselineskip}\subsection*{9.3 The xAct definitions (row 4)}\addcontentsline{toc}{subsection}{9.3 The xAct definitions (row 4)}
\begin{Verbatim}[fontsize=\small,frame=single,framesep=2mm]
CD = LC[g]; RiemannCD = Riemann[CD]; RicciCD = Ricci[CD];
RicciScalarCD = RicciScalar[CD]; EinsteinCD = Einstein[CD];
\end{Verbatim}
{\footnotesize
\begin{longtable}{>{\raggedright\arraybackslash}p{\dimexpr 0.2400\linewidth-2\tabcolsep\relax}>{\raggedright\arraybackslash}p{\dimexpr 0.3600\linewidth-2\tabcolsep\relax}>{\raggedright\arraybackslash}p{\dimexpr 0.4000\linewidth-2\tabcolsep\relax}}
\toprule
\textbf{symbol or term} & \textbf{definition} & \textbf{meaning} \\
\midrule\endfirsthead
\toprule
\textbf{symbol or term} & \textbf{definition} & \textbf{meaning} \\
\midrule\endhead
\endfoot
\bottomrule\endlastfoot
xAct, xCoba, xTras & packages for the Wolfram Language & abstract and component tensor calculus \\
\texttt{g} & a component metric $g_{\mu\nu}$ & the metric whose curvature is computed \\
\texttt{CD = LC[g]} & the Levi-Civita connection of \texttt{g} & after it, the covariant derivative \texttt{CD} is that of $g$ \\
\texttt{RiemannCD}, \texttt{RicciCD} & (9.2) & the Riemann and Ricci tensors of $g$ \\
\texttt{RicciScalarCD}, \texttt{EinsteinCD} & (9.2) & the scalar curvature and the Einstein tensor of $g$ \\
$\{{}^{\lambda}_{\mu\nu}\}$ & (9.2) & Christoffel symbols \\*
\end{longtable}}
\begin{align*}
\{{}^{\lambda}_{\mu\nu}\}&=\tfrac12\,g^{\lambda\rho}\big(\partial_\mu g_{\rho\nu}+\partial_\nu g_{\rho\mu}-\partial_\rho g_{\mu\nu}\big),\qquad R^{\rho}{}_{\sigma\mu\nu}=\partial_\mu\{{}^{\rho}_{\sigma\nu}\}-\partial_\nu\{{}^{\rho}_{\sigma\mu}\}+\{{}^{\rho}_{\mu\lambda}\}\{{}^{\lambda}_{\sigma\nu}\}-\{{}^{\rho}_{\nu\lambda}\}\{{}^{\lambda}_{\sigma\mu}\}, \\
R_{\sigma\nu}&=R^\rho{}_{\sigma\rho\nu},\qquad R=g^{\sigma\nu}R_{\sigma\nu},\qquad G_{\mu\nu}=R_{\mu\nu}-\tfrac12g_{\mu\nu}R .&&(9.2)
\end{align*}
These are the conventions of 7+1.nb (in the Landau–Lifshitz convention for the Riemann tensor); $g^{\lambda\rho}$ is the inverse metric and $\partial_\mu=\partial/\partial x_\mu$.

\Needspace{8\baselineskip}\subsection*{9.4 The mixing angle in matter (row 5)}\addcontentsline{toc}{subsection}{9.4 The mixing angle in matter (row 5)}
\begin{equation*}
\frac{\Delta m^2}{4E}\begin{pmatrix}-\cos2\theta&\sin2\theta\\\sin2\theta&\cos2\theta\end{pmatrix}+\mathrm{diag}(V,0)\quad\Longrightarrow\quad \tan2\theta_m=\frac{\sin2\theta}{\cos2\theta-2EV/\Delta m^2} .
\end{equation*}
{\footnotesize
\begin{longtable}{>{\raggedright\arraybackslash}p{\dimexpr 0.2400\linewidth-2\tabcolsep\relax}>{\raggedright\arraybackslash}p{\dimexpr 0.3600\linewidth-2\tabcolsep\relax}>{\raggedright\arraybackslash}p{\dimexpr 0.4000\linewidth-2\tabcolsep\relax}}
\toprule
\textbf{symbol or term} & \textbf{definition} & \textbf{meaning} \\
\midrule\endfirsthead
\toprule
\textbf{symbol or term} & \textbf{definition} & \textbf{meaning} \\
\midrule\endhead
\endfoot
\bottomrule\endlastfoot
$\Delta m^2$ & $m_2^2-m_1^2$ & difference of the squared masses of two neutrinos \\
$E$ & energy & the energy of the neutrino \\
$\theta$ & the mixing angle in vacuum &  \\
$V$ & $\sqrt2G_Fn_e$ & the potential of the electron neutrino in matter of electron density $n_e$ \\
$\theta_m$ & (row 5) & the mixing angle in matter \\
$\mathrm{diag}(V,0)$ & $\left(\begin{smallmatrix}V&0\\0&0\end{smallmatrix}\right)$ & the matter term of the Hamiltonian \\*
\end{longtable}}
\Needspace{8\baselineskip}\subsection*{9.5 The mode equation of pair creation (row 6)}\addcontentsline{toc}{subsection}{9.5 The mode equation of pair creation (row 6)}
\begin{equation*}
i\,\frac{dc}{d\tau}=\omega\,(\cos\theta\,\sigma_1+\sin\theta\,\sigma_2)\,c,\qquad \omega=K\sqrt{\tau^{-2p_1}+\tau^{-2q_1}},\qquad \tan\theta=\tau^{p_1-q_1} .
\end{equation*}
{\footnotesize
\begin{longtable}{>{\raggedright\arraybackslash}p{\dimexpr 0.2400\linewidth-2\tabcolsep\relax}>{\raggedright\arraybackslash}p{\dimexpr 0.3600\linewidth-2\tabcolsep\relax}>{\raggedright\arraybackslash}p{\dimexpr 0.4000\linewidth-2\tabcolsep\relax}}
\toprule
\textbf{symbol or term} & \textbf{definition} & \textbf{meaning} \\
\midrule\endfirsthead
\toprule
\textbf{symbol or term} & \textbf{definition} & \textbf{meaning} \\
\midrule\endhead
\endfoot
\bottomrule\endlastfoot
$c(\tau)$ & a column of two complex numbers & the amplitude of a mode \\
$\tau$ & $x_0-x_s$ & time since the singularity at $x_0=x_s$ \\
$K$ & a positive number & the scaling variable of the mode (not a multi-index here) \\
$\omega$ & $K\sqrt{\tau^{-2p_1}+\tau^{-2q_1}}$ & the frequency of the mode \\
$\theta$ & $\tan\theta=\tau^{p_1-q_1}$ & the angle of the mode Hamiltonian (not the angle of §9.4) \\
$\sigma_1$, $\sigma_2$ & (9.3) & Pauli matrices \\
$p_1$, $q_1$, $u_1$ & (9.3) & the Kasner exponents of ordinary and auxiliary space in the earliest regime \\*
\end{longtable}}
\begin{align*}
\sigma_1&=\begin{pmatrix}0&1\\1&0\end{pmatrix},\qquad \sigma_2=\begin{pmatrix}0&-i\\i&0\end{pmatrix},\qquad u_1=-1+\frac1{\sqrt2},\qquad p_1=\frac{1+2\sqrt2}7,\qquad q_1=u_1p_1, \\
3p_1+4q_1&=1,\qquad 3p_1^2+4q_1^2=1\qquad(p_1\approx0.5469,\ q_1\approx-0.1602).&&(9.3)
\end{align*}
\Needspace{8\baselineskip}\subsection*{9.6 The minisuperspace action (rows 7 and 8)}\addcontentsline{toc}{subsection}{9.6 The minisuperspace action (rows 7 and 8)}
\begin{align*}
S&=-\frac1{2\kappa}\int n\,a^3b^4\sum_k\alpha_kL_k\,dx_0-\int nW\,dx_0, \\
\mathrm{EL}_n(na^3b^4L_k)&=-2a^3b^4E^{(k)0}{}_0,\qquad \mathrm{EL}_a(na^3b^4L_k)=-6a^2b^4E^{(k)1}{}_1, \\
\mathrm{EL}_b(na^3b^4L_k)&=-8a^3b^3E^{(k)4}{}_4 .
\end{align*}
{\footnotesize
\begin{longtable}{>{\raggedright\arraybackslash}p{\dimexpr 0.2400\linewidth-2\tabcolsep\relax}>{\raggedright\arraybackslash}p{\dimexpr 0.3600\linewidth-2\tabcolsep\relax}>{\raggedright\arraybackslash}p{\dimexpr 0.4000\linewidth-2\tabcolsep\relax}}
\toprule
\textbf{symbol or term} & \textbf{definition} & \textbf{meaning} \\
\midrule\endfirsthead
\toprule
\textbf{symbol or term} & \textbf{definition} & \textbf{meaning} \\
\midrule\endhead
\endfoot
\bottomrule\endlastfoot
$n$, $a$, $b$ & (9.4) & lapse; scale factors of ordinary and of auxiliary space; functions of $x_0$ \\
$S$ & (row 7) & the action restricted to the metrics (9.4) \\
$\kappa$ & $8\pi G_8=M_8^{-6}$ & the eight-dimensional gravitational coupling \\
$\alpha_k$ & $\alpha_0=2\Lambda$, $\alpha_1=1$ & the Lovelock couplings ($\Lambda$ the cosmological constant) \\
$L_k$ & (9.5) & the Lovelock densities: $L_0=1$, $L_1=R$ \\
$\delta^{\mu_1\cdots\mu_p}_{\nu_1\cdots\nu_p}$ & (9.5) & the generalized Kronecker delta \\
$E^{(k)\mu}{}_\nu$ & (9.5) & the Lovelock tensors: $E^{(0)}=-\delta/2$, $E^{(1)}=G$ (the Einstein tensor) \\
$W(a,b)$ & $a^3b^4\varepsilon$ & the comoving energy of the matter ($\varepsilon$ its energy density) \\
$\mathrm{EL}_q$ & (9.6) & the Euler–Lagrange expression for the variable $q=n,a,b$ \\*
\end{longtable}}
\begin{align*}
ds^2&=n^2dx_0^2-a^2\big(dx_1^2+dx_2^2+dx_3^2\big)-b^2\big(dx_4^2+\dots+dx_7^2\big),&&(9.4) \\
L_k&=\frac1{2^k}\delta^{\alpha_1\beta_1\cdots\alpha_k\beta_k}_{\gamma_1\delta_1\cdots\gamma_k\delta_k}R^{\gamma_1\delta_1}{}_{\alpha_1\beta_1}\cdots R^{\gamma_k\delta_k}{}_{\alpha_k\beta_k}, \\
E^{(k)\mu}{}_\nu&=-\frac1{2^{k+1}}\delta^{\mu\alpha_1\beta_1\cdots\alpha_k\beta_k}_{\nu\gamma_1\delta_1\cdots\gamma_k\delta_k}R^{\gamma_1\delta_1}{}_{\alpha_1\beta_1}\cdots R^{\gamma_k\delta_k}{}_{\alpha_k\beta_k}, \\
\delta^{\mu_1\cdots\mu_p}_{\nu_1\cdots\nu_p}&=\det\big(\delta^{\mu_i}_{\nu_j}\big)_{i,j=1}^p,&&(9.5) \\
\mathrm{EL}_q(F)&=\frac{\partial F}{\partial q}-\frac{d}{dx_0}\frac{\partial F}{\partial\dot q}+\frac{d^2}{dx_0^2}\frac{\partial F}{\partial\ddot q}\qquad(\dot q=dq/dx_0).&&(9.6)
\end{align*}
The Greek indices $\alpha_j,\beta_j,\gamma_j,\delta_j$ of (9.5) are summation indices over $0,\dots,7$, not the couplings $\alpha_k$; $R^{\gamma\delta}{}_{\alpha\beta}$ is the Riemann tensor of (9.2) with two indices raised.

\Needspace{8\baselineskip}\subsection*{9.7 The comoving energy of one massless particle (row 8)}\addcontentsline{toc}{subsection}{9.7 The comoving energy of one massless particle (row 8)}
\begin{equation*}
W=\omega=\sqrt{k_a^2/a^2+k_b^2/b^2},\qquad \partial_aW=-\frac{k_a^2}{a^3\omega},\qquad \partial_bW=-\frac{k_b^2}{b^3\omega} .
\end{equation*}
{\footnotesize
\begin{longtable}{>{\raggedright\arraybackslash}p{\dimexpr 0.2400\linewidth-2\tabcolsep\relax}>{\raggedright\arraybackslash}p{\dimexpr 0.3600\linewidth-2\tabcolsep\relax}>{\raggedright\arraybackslash}p{\dimexpr 0.4000\linewidth-2\tabcolsep\relax}}
\toprule
\textbf{symbol or term} & \textbf{definition} & \textbf{meaning} \\
\midrule\endfirsthead
\toprule
\textbf{symbol or term} & \textbf{definition} & \textbf{meaning} \\
\midrule\endhead
\endfoot
\bottomrule\endlastfoot
$k_a$, $k_b$ & constants & the comoving momenta of the particle along ordinary and along auxiliary space \\
$\omega$ & $\sqrt{k_a^2/a^2+k_b^2/b^2}$ & its energy (not the frequency of §9.5) \\
$\partial_a$, $\partial_b$ & $\partial/\partial a$, $\partial/\partial b$ & partial derivatives with respect to the scale factors \\*
\end{longtable}}
\Needspace{8\baselineskip}\subsection*{9.8 The Kaluza–Klein level and the radion (row 9)}\addcontentsline{toc}{subsection}{9.8 The Kaluza–Klein level and the radion (row 9)}
\begin{equation*}
M_E(\beta)=e^{-2\beta}\sqrt{m^2+\mu_n^2e^{-2\beta}},\qquad g_n=\Big|\frac{dM_E}{d\phi}\Big|_{\beta=0}=\frac{2M_n+\mu_n^2/M_n}{\sqrt{12}\,M_{\rm Pl}} .
\end{equation*}
{\footnotesize
\begin{longtable}{>{\raggedright\arraybackslash}p{\dimexpr 0.2400\linewidth-2\tabcolsep\relax}>{\raggedright\arraybackslash}p{\dimexpr 0.3600\linewidth-2\tabcolsep\relax}>{\raggedright\arraybackslash}p{\dimexpr 0.4000\linewidth-2\tabcolsep\relax}}
\toprule
\textbf{symbol or term} & \textbf{definition} & \textbf{meaning} \\
\midrule\endfirsthead
\toprule
\textbf{symbol or term} & \textbf{definition} & \textbf{meaning} \\
\midrule\endhead
\endfoot
\bottomrule\endlastfoot
$\beta$ & $\ln(b/b_0)$ & the radion: the logarithm of the size $b$ of the auxiliary space ($b_0$ its stabilized value, at the minimum of the stabilizing potential) \\
$\phi$ & $\sqrt{12}\,M_{\rm Pl}\beta$ & the canonical radion \\
$M_{\rm Pl}$ & $2.435\times10^{18}$ GeV & the reduced Planck mass \\
$m$ & a mass & the mass of dirac16complex \\
$\mu_n$ & $2\pi|\vec n|/\ell_\perp$ on the auxiliary torus & the Kaluza–Klein mass of the level $n$ ($\ell_\perp$ the size of the torus, $\vec n\in\mathbb Z^4$) \\
$M_E(\beta)$ & (row 9) & the mass of the level in the Einstein frame \\
$M_n$ & $M_E(0)=\sqrt{m^2+\mu_n^2}$ & the mass of the level today \\
$g_n$ & $|dM_E/d\phi|$ at $\beta=0$ & the coupling of the level to the radion \\*
\end{longtable}}
\begin{equation*}
\frac{dM_E}{d\beta}\Big|_{\beta=0}=-2M_n-\frac{\mu_n^2}{M_n},\qquad \frac{d}{d\phi}=\frac1{\sqrt{12}\,M_{\rm Pl}}\frac d{d\beta}\qquad\Longrightarrow\qquad g_n=\frac{2M_n+\mu_n^2/M_n}{\sqrt{12}\,M_{\rm Pl}} .
\end{equation*}
\Needspace{8\baselineskip}\subsection*{9.9 The five momenta (row 10)}\addcontentsline{toc}{subsection}{9.9 The five momenta (row 10)}
\begin{equation*}
\vec p\in\big\{\pm(2,2,4,0,0,0,0),\ \pm(0,0,0,2,2,0,0),\ \vec 0\big\}\qquad(E=5,5,3,3,1).
\end{equation*}
{\footnotesize
\begin{longtable}{>{\raggedright\arraybackslash}p{\dimexpr 0.2400\linewidth-2\tabcolsep\relax}>{\raggedright\arraybackslash}p{\dimexpr 0.3600\linewidth-2\tabcolsep\relax}>{\raggedright\arraybackslash}p{\dimexpr 0.4000\linewidth-2\tabcolsep\relax}}
\toprule
\textbf{symbol or term} & \textbf{definition} & \textbf{meaning} \\
\midrule\endfirsthead
\toprule
\textbf{symbol or term} & \textbf{definition} & \textbf{meaning} \\
\midrule\endhead
\endfoot
\bottomrule\endlastfoot
$\vec p$ & $(p^1,\dots,p^7)$ & the momentum of a plane wave \\
$E$ & $\sqrt{m^2+|\vec p|^2}$ & its energy; with $m=1$: 5, 5, 3, 3, 1 \\
$\vec x$, $\vec y$ & points of space $(x_1,\dots,x_7)$ & the arguments of the field in the sum \\*
\end{longtable}}
\Needspace{8\baselineskip}\subsection*{9.10 The orthogonal matrix of the spinors (row 11)}\addcontentsline{toc}{subsection}{9.10 The orthogonal matrix of the spinors (row 11)}
\begin{equation*}
U=\frac1{\sqrt{2E}}\big(u_1(\vec p),\dots,u_8(\vec p),v_1(-\vec p),\dots,v_8(-\vec p)\big)\qquad\text{satisfies}\qquad U^{\mathsf T}h(\vec p)\,U=\mathrm{diag}(E1_8,-E1_8).
\end{equation*}
{\footnotesize
\begin{longtable}{>{\raggedright\arraybackslash}p{\dimexpr 0.2400\linewidth-2\tabcolsep\relax}>{\raggedright\arraybackslash}p{\dimexpr 0.3600\linewidth-2\tabcolsep\relax}>{\raggedright\arraybackslash}p{\dimexpr 0.4000\linewidth-2\tabcolsep\relax}}
\toprule
\textbf{symbol or term} & \textbf{definition} & \textbf{meaning} \\
\midrule\endfirsthead
\toprule
\textbf{symbol or term} & \textbf{definition} & \textbf{meaning} \\
\midrule\endhead
\endfoot
\bottomrule\endlastfoot
$e_s$ & the unit columns of length 8, $s=1,\dots,8$ &  \\
$M(\vec p)$ & $m1_8+p^i\tau^O_i$ & the real 8×8 block; $MM^{\mathsf T}=E^21_8$ \\
$\alpha_i$ & $\Gamma^0\Gamma^i$ & Dirac’s $\alpha$ matrices; real symmetric \\
$h(\vec p)$ & $p^i\alpha_i+m\Gamma^0$ & the Dirac Hamiltonian of a plane wave; real symmetric \\
$u_s(\vec p)$ & (9.7) & the real spinors of positive energy \\
$v_s(\vec p)$ & (9.7) & the real spinors of negative energy \\
$U$ & (row 11) & the real orthogonal 16×16 matrix of the spinors \\
$E1_8$ & $E$ times the identity &  \\*
\end{longtable}}
\begin{align*}
u_s(\vec p)&=\frac1{\sqrt E}\begin{pmatrix}E\,e_s\\M(\vec p)^{\mathsf T}e_s\end{pmatrix},\qquad h(\vec p)\,u_s=E\,u_s,\qquad h(-\vec p)\,v_s=-E\,v_s,\qquad u_s^{\mathsf T}u_r=v_s^{\mathsf T}v_r=2E\,\delta_{sr}.&&(9.7)
\end{align*}
\Needspace{8\baselineskip}\subsection*{9.11 Completeness and the Fierz identity (row 14)}\addcontentsline{toc}{subsection}{9.11 Completeness and the Fierz identity (row 14)}
\begin{equation*}
\sum_K(\Theta_K)_{cd}(\Theta^K)_{ab}=16\,\delta_{ad}\delta_{bc},
\end{equation*}
\begin{equation*}
\sum_{K\in r}(a\Theta_Kb)(c\Theta^Kd)=\sum_JC_{rJ}\,(a\Theta_Jd)(c\Theta^Jb),\qquad C_{rJ}=\frac1{256}\sum_{K\in r}\mathrm{tr}\big(\Theta_K\Theta_J\Theta^K\Theta^J\big).
\end{equation*}
{\footnotesize
\begin{longtable}{>{\raggedright\arraybackslash}p{\dimexpr 0.2400\linewidth-2\tabcolsep\relax}>{\raggedright\arraybackslash}p{\dimexpr 0.3600\linewidth-2\tabcolsep\relax}>{\raggedright\arraybackslash}p{\dimexpr 0.4000\linewidth-2\tabcolsep\relax}}
\toprule
\textbf{symbol or term} & \textbf{definition} & \textbf{meaning} \\
\midrule\endfirsthead
\toprule
\textbf{symbol or term} & \textbf{definition} & \textbf{meaning} \\
\midrule\endhead
\endfoot
\bottomrule\endlastfoot
$(\Theta_K)_{cd}$ & the entry in row $c$, column $d$ & $a,b,c,d=1,\dots,16$ as subscripts are spinor indices \\
$a,b,c,d$ & commuting columns of 16 numbers & spinors, in the second display \\
$(aMb)$ & $\sum_{i,j}a_iM_{ij}b_j$ & the bilinear of a 16×16 matrix $M$ \\
$K\in r$ & the multi-indices of rank $r$ &  \\
$J$ & a multi-index (any of the 256) &  \\
$C_{rJ}$ & (row 14) & the Fierz coefficients \\
$\mathrm{tr}$ & the trace, the sum of the diagonal entries &  \\*
\end{longtable}}
\Needspace{8\baselineskip}\subsection*{9.12 The notation of sections 7–10 of 7+1.nb (row 15)}\addcontentsline{toc}{subsection}{9.12 The notation of sections 7–10 of 7+1.nb (row 15)}
\begin{gather*}
\Gamma_A=\Gamma^{(1,7)}_A,\qquad \Gamma^{A}=\eta_{(1,7)}^{AA}\Gamma_A,\qquad \Gamma^0=\Gamma_0,\qquad \chi=\chi^{(1,7)}=i\Gamma^{(1,7)}_8,\\ \bar\Psi=\Psi^\dagger\Gamma^0,\qquad \vec x=(x_1,\dots,x_7),\qquad \vec p=(p^1,\dots,p^7).
\end{gather*}
{\footnotesize
\begin{longtable}{>{\raggedright\arraybackslash}p{\dimexpr 0.2400\linewidth-2\tabcolsep\relax}>{\raggedright\arraybackslash}p{\dimexpr 0.3600\linewidth-2\tabcolsep\relax}>{\raggedright\arraybackslash}p{\dimexpr 0.4000\linewidth-2\tabcolsep\relax}}
\toprule
\textbf{symbol or term} & \textbf{definition} & \textbf{meaning} \\
\midrule\endfirsthead
\toprule
\textbf{symbol or term} & \textbf{definition} & \textbf{meaning} \\
\midrule\endhead
\endfoot
\bottomrule\endlastfoot
$\Psi$ & a column of 16 anticommuting operators & the quantized field dirac16complex \\
$\Psi^\dagger$ & its conjugate transpose &  \\
$\bar\Psi$ & $\Psi^\dagger\Gamma^0$ & the Dirac adjoint \\
$\eta_{(1,7)}$ & $\mathrm{diag}(+1,-1,\dots,-1)$ & the same matrix as $\eta$ \\
$\Gamma^{(1,7)}_8$ & (9.8) & the volume element, $(\Gamma^{(1,7)}_8)^2=-1_{16}$ \\
$\chi=\chi^{(1,7)}$ & $i\Gamma^{(1,7)}_8$ & the (complex) chirality matrix, $\chi^2=1_{16}$ \\
$\vec x$, $\vec p$ & $(x_1,\dots,x_7)$, $(p^1,\dots,p^7)$ & space point and momentum \\*
\end{longtable}}
\begin{align*}
\Gamma^{(1,7)}_8&=\Gamma^{(1,7)}_0\Gamma^{(1,7)}_1\cdots\Gamma^{(1,7)}_7=\begin{pmatrix}0&-1_8\\1_8&0\end{pmatrix}.&&(9.8)
\end{align*}
\Needspace{8\baselineskip}\subsection*{9.13 The electroweak scale (row 16)}\addcontentsline{toc}{subsection}{9.13 The electroweak scale (row 16)}
\begin{equation*}
v=(\sqrt2G_F)^{-1/2}=246.2196402\ \text{GeV},\qquad m_W=\frac{gv}2,\qquad m_Z=\frac{\sqrt{g^2+g'^2}\,v}2,\qquad m_H^2=2\lambda v^2 .
\end{equation*}
{\footnotesize
\begin{longtable}{>{\raggedright\arraybackslash}p{\dimexpr 0.2400\linewidth-2\tabcolsep\relax}>{\raggedright\arraybackslash}p{\dimexpr 0.3600\linewidth-2\tabcolsep\relax}>{\raggedright\arraybackslash}p{\dimexpr 0.4000\linewidth-2\tabcolsep\relax}}
\toprule
\textbf{symbol or term} & \textbf{definition} & \textbf{meaning} \\
\midrule\endfirsthead
\toprule
\textbf{symbol or term} & \textbf{definition} & \textbf{meaning} \\
\midrule\endhead
\endfoot
\bottomrule\endlastfoot
$G_F$ & $1.1663788\times10^{-5}$ GeV$^{-2}$ & the Fermi constant \\
$v$ & $(\sqrt2G_F)^{-1/2}=246.2196402$ GeV & the electroweak scale (the vacuum value of the Higgs field) \\
$g$, $g'$ & dimensionless & the gauge couplings of $SU(2)_L$ and $U(1)_Y$ \\
$\lambda$ & dimensionless & the quartic coupling of the Higgs field \\
$m_W$, $m_Z$, $m_H$ & (row 16) & the masses of the W, Z and Higgs bosons at tree level (lowest order) \\*
\end{longtable}}
\Needspace{8\baselineskip}\subsection*{9.14 The matrix notation (row 17)}\addcontentsline{toc}{subsection}{9.14 The matrix notation (row 17)}
\begin{gather*}
1_n=(\delta_{ij})_{i,j=1}^n,\qquad (M^\dagger)_{ij}=\overline{M_{ji}},\qquad (M^*)_{ij}=\overline{M_{ij}},\\ (A\otimes B)_{(i-1)p+k,\,(j-1)q+l}=A_{ij}B_{kl}\ \ (B\ \text{of size}\ p\times q).
\end{gather*}
{\footnotesize
\begin{longtable}{>{\raggedright\arraybackslash}p{\dimexpr 0.2400\linewidth-2\tabcolsep\relax}>{\raggedright\arraybackslash}p{\dimexpr 0.3600\linewidth-2\tabcolsep\relax}>{\raggedright\arraybackslash}p{\dimexpr 0.4000\linewidth-2\tabcolsep\relax}}
\toprule
\textbf{symbol or term} & \textbf{definition} & \textbf{meaning} \\
\midrule\endfirsthead
\toprule
\textbf{symbol or term} & \textbf{definition} & \textbf{meaning} \\
\midrule\endhead
\endfoot
\bottomrule\endlastfoot
$1_n$ & $(\delta_{ij})$ & the $n\times n$ identity \\
$M^{\mathsf T}$ & $(M^{\mathsf T})_{ij}=M_{ji}$ & the transpose \\
$M^\dagger$ & $(M^\dagger)_{ij}=\overline{M_{ji}}$ & the conjugate transpose \\
$M^*$ & $(M^*)_{ij}=\overline{M_{ij}}$ & the complex conjugate \\
$\overline z$ & the complex conjugate of the number $z$ &  \\
$A\otimes B$ & (row 17) & the Kronecker product; in Mathematica \texttt{KroneckerProduct[A, B]} \\*
\end{longtable}}
\Needspace{12\baselineskip}\section*{10. File and program names}\addcontentsline{toc}{section}{10. File and program names}
{\footnotesize
\begin{longtable}{>{\raggedright\arraybackslash}p{\dimexpr 0.4500\linewidth-2\tabcolsep\relax}>{\raggedright\arraybackslash}p{\dimexpr 0.5500\linewidth-2\tabcolsep\relax}}
\toprule
\textbf{name} & \textbf{what it is} \\
\midrule\endfirsthead
\toprule
\textbf{name} & \textbf{what it is} \\
\midrule\endhead
\endfoot
\bottomrule\endlastfoot
\texttt{7plus1\_\allowbreak{}status\_\allowbreak{}2026-10-28.xlsx} & the spreadsheet: the sheets Summary, Commits, Documents, Fixes, Checks, Private files \\
\texttt{tools/\allowbreak{}status\_\allowbreak{}2026-10-28.json} & the data file: every number of the spreadsheet \\
\texttt{tools/\allowbreak{}make\_\allowbreak{}status\_\allowbreak{}xlsx.py} & writes the spreadsheet from the data file \\
\texttt{tools/\allowbreak{}make\_\allowbreak{}docs.py} & writes every document (Markdown, LaTeX, PDF) and runs the twelve check programs \\
\texttt{tools/\allowbreak{}doc\_\allowbreak{}status.py}, \texttt{tools/\allowbreak{}doc\_\allowbreak{}status\_\allowbreak{}symbols.py}, \texttt{tools/\allowbreak{}doc\_\allowbreak{}status\_\allowbreak{}provenance.py}, \texttt{tools/\allowbreak{}doc\_\allowbreak{}status\_\allowbreak{}common.py} & write the document on the state of the repository, this dictionary and the provenance record \\
\texttt{build/\allowbreak{}docs/\allowbreak{}runs.json} & the outputs of the twelve programs in the last build (build/ is not published) \\
\texttt{tools/\allowbreak{}eval\_\allowbreak{}log\_\allowbreak{}7p1.txt} & the record of the evaluation of 7+1.nb \\
\texttt{tools/\allowbreak{}verify\_\allowbreak{}7p1\_\allowbreak{}nb.py} & parses the evaluated 7+1.nb and compares it with the Python programs and the backups \\
\texttt{tools/\allowbreak{}ref44.py}, \texttt{tools/\allowbreak{}ref71.py}, \texttt{tools/\allowbreak{}extra71.py} & rebuild 4+4.nb and the objects of 7+1.nb; further identities \\
\texttt{tools/\allowbreak{}ref\_\allowbreak{}dirac16.py}, \texttt{tools/\allowbreak{}ref\_\allowbreak{}gravity.py}, \texttt{tools/\allowbreak{}ref\_\allowbreak{}ks\_\allowbreak{}gravity.py}, \texttt{tools/\allowbreak{}ref\_\allowbreak{}sm.py} & independent checks of sections 7–9, 11–14, 15–16 and 17–19 of 7+1.nb \\
\texttt{tools/\allowbreak{}ref\_\allowbreak{}coupled.py}, \texttt{tools/\allowbreak{}ref\_\allowbreak{}dark.py}, \texttt{tools/\allowbreak{}ref\_\allowbreak{}torus\_\allowbreak{}brane.py}, \texttt{tools/\allowbreak{}ref\_\allowbreak{}radion.py} & independent checks of sections 20–23, 24, 25 and 26 of 7+1.nb \\
\texttt{tools/\allowbreak{}doc\_\allowbreak{}stage3\_\allowbreak{}common.py} & shared data of the documents of sections 11–26; its SHA-256 is recorded in the provenance files \\
7+1.nb & the notebook of the project (Wolfram Language 15.0.1), evaluated \\
Einstein-Lovelock-7+1-Nash.nb & the author’s reference notebook for the Lovelock tensors with xAct \\*
\end{longtable}}
\Needspace{12\baselineskip}\section*{11. Symbols that look alike}\addcontentsline{toc}{section}{11. Symbols that look alike}
{\footnotesize
\begin{longtable}{>{\raggedright\arraybackslash}p{\dimexpr 0.4200\linewidth-2\tabcolsep\relax}>{\raggedright\arraybackslash}p{\dimexpr 0.5800\linewidth-2\tabcolsep\relax}}
\toprule
\textbf{symbol} & \textbf{is not} \\
\midrule\endfirsthead
\toprule
\textbf{symbol} & \textbf{is not} \\
\midrule\endhead
\endfoot
\bottomrule\endlastfoot
$\Gamma_A$, $\Gamma^A$ (Dirac matrices, one index) & $\{{}^{\lambda}_{\mu\nu}\}$ (Christoffel symbols, never written $\Gamma$ here) \\
$K$ (multi-index, §9.1, §9.11) & $K$ (the scaling variable of a mode, §9.5) \\
$\theta$ (mixing angle in vacuum, §9.4) & $\theta$ (angle of the mode Hamiltonian, §9.5); $\theta_k$, $\theta_K$ (Grassmann generators, §9.2); $\theta_m$ (mixing angle in matter, §9.4) \\
$\sigma_r$ (0 or 1, §9.1) & $\sigma_1$, $\sigma_2$ (Pauli matrices, §9.5) \\
$\omega$ (frequency of a mode, §9.5) & $\omega$ (energy of one massless particle, §9.7) \\
$E$ (energy) & $E^{(k)\mu}{}_\nu$ (Lovelock tensors, §9.6); $E1_8$ ($E$ times the identity) \\
$a$, $b$ (scale factors, §9.6) & $a,b,c,d$ (spinors or spinor indices, §9.11) \\
$c(\tau)$ (mode amplitude, §9.5) & $c$ (a spinor or a spinor index, §9.11); $c$ (a cell of the spreadsheet, section 6) \\
$\beta$ (radion, §9.8) & $\alpha_k$ (Lovelock couplings) and the summation indices $\alpha_j,\beta_j$ of (9.5) \\
$g$ (a metric, §9.3) & $g$, $g'$ (gauge couplings, §9.13); $g_n$ (coupling of a level to the radion, §9.8) \\
$R$ (scalar curvature, §9.3) & $R$ (lines removed, section 2); $R$ (a range of cells, section 6) \\
$s$ (stretchability, section 3) & $s$ (spinor label $s=1,\dots,8$, §9.10) \\
$\ell$ (a line of output, section 5) & $\ell_\perp$ (the size of the auxiliary torus, §9.8) \\
$\eta$, $\eta_{(1,7)}$ (the same matrix) & $\eta_{71}=-\eta$ (the sign convention of 4+4.nb, section 1) \\*
\end{longtable}}
\end{document}
